
\documentclass{article}
\usepackage{colortbl}
\usepackage{makecell}
\usepackage{multirow}
\usepackage{supertabular}

\begin{document}

\newcounter{utterance}

\centering \large Interaction Transcript for game `cladder', experiment `full\_v1.5\_default', episode 2340 with Qwen3.5{-}9B.
\vspace{24pt}

{ \footnotesize  \setcounter{utterance}{1}
\setlength{\tabcolsep}{0pt}
\begin{supertabular}{c@{$\;$}|p{.15\linewidth}@{}p{.15\linewidth}p{.15\linewidth}p{.15\linewidth}p{.15\linewidth}p{.15\linewidth}}
    \# & \multicolumn{2}{c}{Player} && \multicolumn{2}{c}{Game Master} \\
    \hline

    \theutterance \stepcounter{utterance}  
    & & & \multicolumn{4}{p{0.6\linewidth}}{
        \cellcolor[rgb]{0.9,0.9,0.9}{
            \makecell[{{p{\linewidth}}}]{
                \texttt{\tiny{[P1$\langle$GM]}}
                \texttt{You are an expert in causal inference. The following question is not a typical commonsense query, but rather a meticulously designed question created by a professor specializing in causal inference, intended to assess the students' mastery of the course content.} \\
\texttt{Imagine a self{-}contained, hypothetical world with only the following conditions, and without any unmentioned factors or causal relationships: Eating citrus has a direct effect on vitmain C. Vitmain C has a direct effect on scurvy. For patients not consuming citrus, the probability of scurvy is 73\%. For patients consuming citrus, the probability of scurvy is 54\%. Does eating citrus positively affect scurvy through vitmain C?} \\
\texttt{Start your answer with 'Yes' or 'No', followed by additional reasoning or evidence to support your explanation.} \\
            }
        }
    }
     \\ \\

    \theutterance \stepcounter{utterance}  
    & \multicolumn{4}{p{0.6\linewidth}}{
        \cellcolor[rgb]{0.561,0.737,0.561}{
            \makecell[{{p{\linewidth}}}]{
                \texttt{\tiny{[P1$\rangle$GM]}}
                \texttt{No} \\
\\ 
\texttt{**Reasoning:**} \\
\\ 
\texttt{To determine if eating citrus *positively* affects scurvy, we must analyze the direction of the causal effect described in the hypothetical world.} \\
\\ 
\texttt{1.  **Analyze the Causal Chain:**} \\
\texttt{    *   The prompt states: "Eating citrus has a direct effect on vitamin C."} \\
\texttt{    *   The prompt states: "Vitamin C has a direct effect on scurvy."} \\
\texttt{    *   In the context of disease and nutrition, the presence of Vitamin C *prevents* or *reduces* the occurrence of scurvy. Therefore, the causal relationship between Vitamin C and Scurvy is **negative** (inverse). An increase in Vitamin C leads to a decrease in Scurvy.} \\
\\ 
\texttt{2.  **Analyze the Empirical Data:**} \\
\texttt{    *   Probability of Scurvy (No Citrus) = 73\%} \\
\texttt{    *   Probability of Scurvy (Citrus) = 54\%} \\
\texttt{    *   The difference (Risk Difference) = \$54\textbackslash{}\% {-} 73\textbackslash{}\% = {-}19\textbackslash{}\%\$.} \\
\texttt{    *   This indicates that eating citrus **reduces** the probability of developing scurvy.} \\
\\ 
\texttt{3.  **Evaluate the Question's Premise:**} \\
\texttt{    *   The question asks: "Does eating citrus **positively** affect scurvy...?"} \\
\texttt{    *   In causal inference, a "} \\
            }
        }
    }
    & & \\ \\

    \theutterance \stepcounter{utterance}  
    & & & \multicolumn{2}{p{0.3\linewidth}}{
        \cellcolor[rgb]{0.9,0.9,0.9}{
            \makecell[{{p{\linewidth}}}]{
                \texttt{\tiny{[GM$|$GM]}}
                \texttt{No} \\
            }
        }
    }
    & & \\ \\

    \theutterance \stepcounter{utterance}  
    & & & \multicolumn{2}{p{0.3\linewidth}}{
        \cellcolor[rgb]{0.9,0.9,0.9}{
            \makecell[{{p{\linewidth}}}]{
                \texttt{\tiny{[GM$|$GM]}}
                \texttt{no} \\
            }
        }
    }
    & & \\ \\

    \theutterance \stepcounter{utterance}  
    & & & \multicolumn{2}{p{0.3\linewidth}}{
        \cellcolor[rgb]{0.9,0.9,0.9}{
            \makecell[{{p{\linewidth}}}]{
                \texttt{\tiny{[GM$|$GM]}}
                \texttt{game\_result = WIN} \\
            }
        }
    }
    & & \\ \\

\end{supertabular}
}

\end{document}
